% ======================================================================
%  sec16.tex —— 第 16 节 凝聚模
%  底本：原书 pp.78–82（印刷页 = PDF 页）
%  原书标题：16. Coherent modules
% ======================================================================

\section{凝聚模}\label{sec:16}

\subsection{有限型模}\label{sec:16.1}

\begin{definition}\label{def:16.1.1}
设 $(X,\cO_X)$ 是环化空间，$\cF\in\omod$。若存在截面 $s_i\in\cF(X)$（$i\in I$），使态射
\[
\bigoplus_{i\in I}\cO_X \xrightarrow{(s_i)_{i\in I}} \cF
\]
在阿贝尔范畴 $\omod$ 中是满射，则称 $\cF$ 由截面 ${(s_i)}_{i\in I}$ 生成。
\end{definition}

\begin{definition}\label{def:16.1.2}
设 $(X,\cO_X)$ 是环化空间。若对任意点 $x\in X$，存在 $x$ 的开邻域 $U\subseteq X$，使 $\cF|_U$ 由有限多个截面生成，则称 $\cF\in\omod$ 为\textbf{有限型}（of finite type）。
\end{definition}

\begin{lemma}\label{lem:16.1.3}
设 $f\colon (X,\cO_X)\to(Y,\cO_Y)$ 是环化空间之间的态射。则 $f^*$ 把有限型 $\cO_Y$-模送到有限型 $\cO_X$-模。
\end{lemma}

\begin{proof}
由 $f^*$ 右正合这一事实即得。
\end{proof}

\begin{lemma}\label{lem:16.1.4}
设 $(X,\cO_X)$ 是环化空间，$0\to\cF_1\xrightarrow{\alpha}\cF_2\xrightarrow{\beta}\cF_3\to 0$ 是 $\cO_X$-模的短正合列。
\begin{enumerate}[label=(\arabic*),leftmargin=2em,itemsep=0.25ex]
\item 若 $\cF_2$ 是有限型，则 $\cF_3$ 亦然。
\item 若 $\cF_1$ 与 $\cF_3$ 是有限型，则 $\cF_2$ 亦然。
\end{enumerate}
\end{lemma}

\begin{proof}
(1) 显然。

(2) 设 $x\in X$，取 $x$ 的开邻域 $U\subseteq X$，使 $\cF_1|_U$ 与 $\cF_3|_U$ 分别由截面 $a_i\in\cF_1(U)$（$i\in I$）与 $c_j\in\cF_3(U)$（$j\in J$）生成，其中 $|I|,|J|<\infty$。因 $\cF_2\to\cF_3$ 是满射，对每个 $j\in J$ 存在 $x$ 的开邻域 $U_j\subseteq U$ 与截面 $b_j\in\cF_2(U_j)$，使 $\beta(b_j)=c_j|_{U_j}$。置 $V:=\bigcap_{j\in J}U_j$。容易看出，$\cF_2|_V$ 由 $\alpha(a_i)|_V$（$i\in I$）与 $b_j|_V$（$j\in J$）生成。
\end{proof}

\begin{warning}\label{warn:16.1.5}
有限型的 $\cO_X$-模构成的范畴未必是阿贝尔范畴。
\end{warning}

\begin{yizhu}
警告的成因在核：有限型模之间的态射，其核未必有限型，于是这类模对核不封闭（对余核则封闭，见引理 \ref{lem:16.1.4}(1)）。凝聚模的定义（\ref{def:16.4.1}）正是补上这一条：任意有限个截面的核仍要求有限型。补上之后，$\mathit{Coh}(\cO_X)$ 对任意环化空间都是 $\mathit{Mod}(\cO_X)$ 的阿贝尔子范畴。
[来源：Stacks Project, Tag 01BU（Section 17.12, Coherent modules 的开篇），\url{https://stacks.math.columbia.edu/tag/01BU}（2026-11-17 访问）。]
\end{yizhu}

\begin{lemma}\label{lem:16.1.6}
设 $(X,\cO_X)$ 是环化空间，$\cF$ 是有限型 $\cO_X$-模。则 $\mathrm{supp}(\cF):=\{x\in X\mid \cF_x\neq 0\}$ 是 $X$ 的闭子集。
\end{lemma}

\begin{proof}
只需说明：若某点 $x\in X$ 处 $\cF_x=0$，则存在 $x$ 的开邻域 $V\subseteq X$ 使 $\cF|_V\simeq 0$。取 $U$ 使 $\cF|_U$ 由 $s_i\in\cF(U)$（$i\in I$，$|I|<\infty$）生成。由 $\cF_x=0$ 得 $(s_i)_x=0$；故对每个 $i\in I$，存在 $x$ 的开邻域 $U_i\subseteq U$ 使 $s_i|_{U_i}=0$。置 $V:=\bigcap_{i\in I}U_i$。由于 $\cF|_U$ 由 $s_i$ 生成，$\cO_V$-模 $\cF|_V$ 由 $s_i|_V$（$i\in I$）生成；而 $s_i|_V=0$，故 $\cF|_V=0$，如所欲证。
\end{proof}

\begin{lemma}\label{lem:16.1.7}
设 $X$ 是概形，$\cF$ 是有限型的拟凝聚 $\cO_X$-模。则对任意仿射开子概形 $U\subseteq X$，$\cF(U)$ 是有限生成 $\cO_X(U)$-模。
\end{lemma}

\begin{proof}
可设 $X=U=\Spec(A)$ 是仿射概形。此时命题化归为下面的交换代数习题 \ref{exc:16.1.8}。
\end{proof}

\begin{yizhu}
「有限型」与「拟凝聚」是两条独立的约束：前者只管局部生成元的个数，后者管层能否由仿射开集上的模黏合出来。在仿射概形上二者交叠成一句交换代数：$\widetilde{M}$ 有限型当且仅当 $M$ 是有限 $A$-模，这正是本引理与其逆。警告 \ref{warn:16.2.4} 说明单有有限型推不出拟凝聚，故本引理必须同时假设两者。
[来源：Stacks Project, Tag 01PB（Lemma 28.17.1，$\widetilde{M}$ 有限型当且仅当 $M$ 有限），\url{https://stacks.math.columbia.edu/tag/01PB}（2026-11-17 访问）。]
\end{yizhu}

\begin{exercise}\label{exc:16.1.8}
设 $A$ 是交换环，$f_1,\dots,f_n\in A$ 使 $(f_1,\dots,f_n)=A$。若 $A$-模 $M$ 满足每个 $M_{f_i}$ 是有限生成 $A_{f_i}$-模，则 $M$ 是有限生成 $A$-模。
\end{exercise}

\begin{corollary}\label{cor:16.1.9}
设 $X$ 是概形，$\cF$ 是有限型的拟凝聚 $\cO_X$-模。则对任意 $x\in X$，$\cF_x$ 是有限生成 $\cO_{X,x}$-模。
\end{corollary}

利用 Nakayama 引理，得到下述结果。

\begin{corollary}\label{cor:16.1.10}
设 $X$ 是概形，$\cF$ 是有限型的拟凝聚 $\cO_X$-模。则 $x\in\mathrm{supp}(\cF)$ 当且仅当 $\cF$ 在 $x$ 处的纤维 $\cF_x\otimes_{\cO_{X,x}}\kappa_x$ 非零。
\end{corollary}

\begin{construction}\label{con:16.1.11}
设 $X$ 是概形，$\cF$ 是有限型的拟凝聚 $\cO_X$-模。对任意仿射开子概形 $U\subseteq X$，令
\[
\mathrm{Ann}_{\cO_X(U)}(\cF(U))\subseteq\cO_X(U)
\]
为 $\cF(U)$ 的零化子。由引理 \ref{lem:16.1.7}，$\cF(U)$ 是有限生成 $\cO_X(U)$-模。由此对任意 $f\in\cO_X(U)$ 有
\[
\mathrm{Ann}_{\cO_X(U)}(\cF(U))_f \simeq \mathrm{Ann}_{\cO_X(U_f)}(\cF(U_f)),
\]
其中 $U_f:=U\cap D(f)$。这说明存在唯一的拟凝聚理想
\[
\mathrm{Ann}_{\cO_X}(\cF)\subseteq\cO_X,
\]
使对任意仿射开子概形 $U\subseteq X$ 有
\[
(\mathrm{Ann}_{\cO_X}(\cF))(U)=\mathrm{Ann}_{\cO_X(U)}(\cF(U)).
\]
称 $\mathrm{Ann}_{\cO_X}(\cF)$ 为 $\cF$ 的\textbf{零化子}（annihilator）。
\end{construction}

\begin{proposition}\label{prop:16.1.12}
设 $X$ 是概形，$\cF$ 是有限型的拟凝聚 $\cO_X$-模，并设 $i\colon Y\to X$ 是对应于 $\mathrm{Ann}_{\cO_X}(\cF)$ 的闭浸入。则存在本质唯一的拟凝聚 $\cO_Y$-模 $\cM$，连同同构 $\alpha\colon\cF\xrightarrow{\simeq}i_*\cM$。此外还有：
\begin{enumerate}[label=(\arabic*),leftmargin=2em,itemsep=0.25ex]
\item 拟凝聚模 $\cM$ 是有限型。
\item $\mathrm{supp}(\cM)=Y$，从而 $\mathrm{supp}(\cF)=i(Y)$。
\item 对任意闭浸入 $i'\colon Y'\to X$ 与 $\cM'\in\mathsf{QCoh}(Y')$，连同同构 $\alpha'\colon\cF\xrightarrow{\simeq}(i')_*\cM'$，存在唯一的 $X$-态射 $f\colon Y\to Y'$ 与唯一的同构 $\beta\colon\cM'\xrightarrow{\simeq}f_*\cM$，使 $(i')_*(\beta)\circ\alpha'=\alpha$。
\end{enumerate}
称 $Y$ 为 $\cF$ 的\textbf{概形论支集}（scheme theoretic support）。
\end{proposition}

\begin{proof}
对任意如 (3) 中的三元组 $(i',\cM',\alpha')$，有典范同构
\begin{equation}
(i')^*(\cF)\xrightarrow{(i')^*(\alpha)}(i')^*\circ(i')_*(\cM')\xrightarrow{\simeq}\cM'. \tag{16.1}
\end{equation}
于是 $(i',\cM',\alpha')$ 若存在则本质唯一；并且必须有 $\cM\simeq i^*\cF$，而 $\alpha$ 由典范态射
\begin{equation}
\cF\to i_*\circ i^*(\cF) \tag{16.2}
\end{equation}
给出。故要证 $(\cM,\alpha)$ 存在，只需证明 (16.2) 是同构。由构造 \ref{con:16.1.11}，对任意仿射开子概形 $U\subseteq X$ 有
\[
\cF(U)\simeq \cO_X(U)/\mathrm{Ann}_{\cO_X(U)}(\cF(U))\otimes_{\cO_X(U)}\cF(U)\simeq \cO_Y(U\cap Y)\otimes_{\cO_X(U)}\cF(U)\simeq (i_*\circ i^*(\cF))(U).
\]
这说明 (16.2) 确为同构。

断言 (1) 由引理 \ref{lem:16.1.3} 得到。

为证断言 (2)，可设 $Y=\Spec(A)$ 是仿射概形，$\cM=\widetilde{M}$，且 $\mathrm{Ann}_A(M)=0$。利用 $M$ 是有限生成 $A$-模这一事实，容易看出对任意素理想 $\mathfrak{p}\subseteq A$ 有 $M_{\mathfrak{p}}\neq 0$。故 $\mathrm{supp}(\cM)=Y$，如所欲证。

为证断言 (3)，设 $I'$ 是 $i'\colon Y'\to X$ 的定义理想。容易看出
\[
\mathrm{Ann}_{\cO_X}(\cF)\simeq\mathrm{Ann}_{\cO_X}((i')_*\cM')\supseteq I'.
\]
于是由闭子概形的分类（定理 \ref{thm:15.1.1}），得到唯一的 $X$-态射 $f\colon Y\to Y'$。此外有典范同构
\[
\cM'\xrightarrow{\simeq}(i')^*(\cF)\simeq(i')^*\circ i_*(\cM)\simeq(i')^*\circ(i')_*\circ f_*(\cM)\simeq f_*(\cM).
\]
由此得到所求的同构 $\beta$；图追可证其唯一。
\end{proof}

\begin{yizhu}
断言 (2) 的证明省了一步。若 $M_{\mathfrak{p}}=0$，则 $\mathrm{Ann}_A(M)_{\mathfrak{p}}=\mathrm{Ann}_{A_{\mathfrak{p}}}(M_{\mathfrak{p}})=A_{\mathfrak{p}}$，于是存在 $f\notin\mathfrak{p}$ 零化 $M$，与 $\mathrm{Ann}_A(M)=0$ 矛盾。这里「零化子与局部化交换」用到了 $M$ 有限生成：对一般模只有包含关系，右端可能更大，上述论证失效。
[依据：原书 p.80（命题 \ref{prop:16.1.12} 的证明中「容易看出」一步）与 p.79 构造 \ref{con:16.1.11}。]
\end{yizhu}

\begin{yizhu}
原文把这条写成命题-定义：它既断言 $Y$ 与 $\cM$ 的存在唯一，又引入「概形论支集」这个名称。概形论支集比闭集 $\mathrm{supp}(\cF)$ 多记了重数。例如在 $X=\A^1$ 上取 $\cF=\cO_X/(x^2)$：作为集合的支集只有原点，而零化子是 $(x^2)$，故概形论支集为 $\Spec k[x]/(x^2)$。
[依据：原书 p.80 命题 \ref{prop:16.1.12} 及其断言 (2)。]
\end{yizhu}

\subsection{有限表现模}\label{sec:16.2}

\begin{definition}\label{def:16.2.1}
设 $(X,\cO_X)$ 是环化空间。若对任意点 $x\in X$，存在 $x$ 的开邻域 $U\subseteq X$ 与正合列
\[
\bigoplus_{i\in I}\cO_U\to\bigoplus_{j\in J}\cO_U\to\cF|_U\to 0,
\]
其中 $|I|,|J|<\infty$，则称 $\cF\in\omod$ 为\textbf{有限表现}（of finite presentation）。
\end{definition}

下面这个结果是显然的。

\begin{lemma}\label{lem:16.2.2}
设 $(X,\cO_X)$ 是环化空间，$\cF$ 是有限表现的 $\cO_X$-模。则 $\cF$ 是有限型的。
\end{lemma}

\begin{lemma}\label{lem:16.2.3}
设 $X$ 是概形，$\cF$ 是有限表现的 $\cO_X$-模。则 $\cF$ 是拟凝聚的。
\end{lemma}

\begin{proof}
对每个 $x\in X$，取 $x$ 的开邻域 $U$，使有正合列
\[
\bigoplus_{i\in I}\cO_U\to\bigoplus_{j\in J}\cO_U\to\cF|_U\to 0,
\]
其中 $|I|,|J|<\infty$。则 $\cF|_U$ 是拟凝聚 $\cO_U$-模的有限余极限，故拟凝聚。由于拟凝聚性在概形上是局部的，$\cF$ 是拟凝聚的。
\end{proof}

\begin{warning}\label{warn:16.2.4}
设 $X$ 是概形。有限型的 $\cO_X$-模未必拟凝聚。事实上，对任意非拟凝聚理想 $\cI\subseteq\cO_X$，$\cO_X/\cI$ 是有限型的，但不拟凝聚。
\end{warning}

下面的习题是引理 \ref{lem:16.1.4} 的类比。

\begin{exercise}\label{exc:16.2.5}
设 $(X,\cO_X)$ 是环化空间，$0\to\cF_1\xrightarrow{\alpha}\cF_2\xrightarrow{\beta}\cF_3\to 0$ 是 $\cO_X$-模的短正合列。
\begin{enumerate}[label=(\arabic*),leftmargin=2em,itemsep=0.25ex]
\item 若 $\cF_1$ 有限型、$\cF_2$ 有限表现，则 $\cF_3$ 有限表现。
\item 若 $\cF_1$ 与 $\cF_3$ 有限表现，则 $\cF_2$ 亦然。
\item 若 $\cF_2$ 有限型、$\cF_3$ 有限表现，则 $\cF_1$ 有限型。
\end{enumerate}
\end{exercise}

\begin{warning}\label{warn:16.2.6}
有限表现的 $\cO_X$-模构成的范畴未必是阿贝尔范畴。
\end{warning}

\subsection{有限表现模与内 Hom}\label{sec:16.3}

口号是：当 $M$ 有限表现时，内 Hom 函子 $\cHom_{\cO_X}(M,-)$ 表现良好。

\begin{lemma}\label{lem:16.3.1}
设 $(X,\cO_X)$ 是环化空间，$M,N\in\omod$。若 $M$ 有限表现，则典范映射
\[
\cHom_{\cO_X}(M,N)_x\to\Hom_{\cO_{X,x}}(M_x,N_x)
\]
是同构。
\end{lemma}

\begin{proof}
缩小 $X$，可设 $M$ 处在正合列
\[
\bigoplus_{I}\cO_X\to\bigoplus_{J}\cO_X\to M\to 0
\]
中，其中 $I,J$ 是有限集。于是有交换图
\[
\begin{tikzcd}[column sep=small]
0 \arrow[r] & \cHom_{\cO_X}(M,N)_x \arrow[r] \arrow[d] & \cHom_{\cO_X}(\cO_X^{\oplus I},N)_x \arrow[r] \arrow[d] & \cHom_{\cO_X}(\cO_X^{\oplus J},N)_x \arrow[d] \\
0 \arrow[r] & \Hom_{\cO_{X,x}}(M_x,N_x) \arrow[r] & \Hom_{\cO_{X,x}}\bigl(\cO_{X,x}^{\oplus I},N_x\bigr) \arrow[r] & \Hom_{\cO_{X,x}}\bigl(\cO_{X,x}^{\oplus J},N_x\bigr)
\end{tikzcd}
\]
且两行都正合。故只需说明第二个与第三个纵向映射是同构，而这是显然的：这两个映射的来源与目标都可以等同于 $\bigoplus_I N_x$ 或 $\bigoplus_J N_x$。
\end{proof}

\begin{yizhu}
原文在 16.3 里把内 Hom 一律写成 $\Hom_{\cO_X}(M,N)$，与整体截面的 $\Hom$ 同形。译文按层 $\cHom$ 排：带下划线的是 $\omod$ 中的对象，写成 $\Hom$ 的是模同态群。图中两行正合都来自 $\cHom$ 对第一个变量反变且左正合；首列左端的 $0$ 即由此而来。纵向映射之所以好算，是因为 $\cHom(\cO_X^{\oplus I},N)\simeq N^{\oplus I}$，而取茎与有限直和交换。
[依据：原书 p.81 引理 \ref{lem:16.3.1} 的证明与其前的口号。]
\end{yizhu}

类似地可以证明下述结果。

\begin{lemma}\label{lem:16.3.2}
设 $(X,\cO_X)$ 是环化空间，$M$ 是有限表现的 $\cO_X$-模。则函子
\[
\cHom_{\cO_X}(M,-)\colon\omod\to\omod
\]
保持滤过余极限。
\end{lemma}

\begin{lemma}\label{lem:16.3.3}
设 $X$ 是概形，$M$ 是有限表现的 $\cO_X$-模。则函子
\[
\cHom_{\cO_X}(M,-)\colon\omod\to\omod
\]
保持拟凝聚模。所得函子
\[
\cHom_{\cO_X}(M,-)\colon\mathsf{QCoh}(X)\to\mathsf{QCoh}(X)
\]
保持滤过余极限。
\end{lemma}

\begin{yizhu}
「有限表现」在这里的作用是让生成元与关系都有限，从而 $\cHom(M,-)$ 能逐项与余极限交换：把 $M$ 写成 $\cO_X^{\oplus I}\to\cO_X^{\oplus J}\to M\to 0$（$I,J$ 有限），再用 $\cHom(\cO_X^{\oplus I},-)$ 与有限直和、余极限的交换性。若不要求有限表现，这条交换性一般不成立：无限秩自由模上的 $\cHom$ 会把直和变成直积，余极限无法穿过。
[依据：原书 p.81 引理 \ref{lem:16.3.2}、\ref{lem:16.3.3} 与其前的口号。]
\end{yizhu}

\begin{exercise}\label{exc:16.3.4}
设 $X=\Spec(A)$ 是仿射概形，$M,N$ 是 $A$-模。若 $M$ 有限表现，证明 $\widetilde{M}$ 有限表现，且 $\cHom_{\cO_X}(\widetilde{M},\widetilde{N})$ 是 $A$-模 $\Hom_A(M,N)$ 对应的拟凝聚 $\cO_X$-模。
\end{exercise}

\subsection{凝聚模}\label{sec:16.4}

\begin{definition}\label{def:16.4.1}
设 $(X,\cO_X)$ 是环化空间。若 $\cO_X$-模 $\cF$ 是有限型的，且对任意开子集 $U\subseteq X$ 与任意有限个截面 $s_i\in\cF(U)$（$i\in I$），核
\[
\ker\Bigl(\bigoplus_{i\in I}\cO_U\xrightarrow{(s_i)_{i\in I}}\cF|_U\Bigr)
\]
都是有限型的，则称 $\cF$ 为\textbf{凝聚的}（coherent）。
\end{definition}

上述概念一般表现不良。本课程只在局部诺特概形上考虑凝聚模。

\begin{yizhu}
在局部诺特概形上，凝聚这一条可以化归为两个熟悉的条件：$\cF$ 凝聚当且仅当 $\cF$ 拟凝聚且有限型，也当且仅当在每个仿射开集 $U=\Spec(A)$ 上 $\cF\simeq\widetilde{M}$，$M$ 是有限 $A$-模。原文在定义之后立刻声明只在局部诺特概形上考虑凝聚模，原因即在此：一般环化空间上「核有限型」难以检验，连凝聚层的拉回都未必凝聚。
[来源：Stacks Project, Tag 01XZ（Lemma 30.9.1，局部诺特概形上凝聚模的三种刻画），\url{https://stacks.math.columbia.edu/tag/01XZ}（2026-11-17 访问）。]
\end{yizhu}
